% UC04 -- Manter humor

\begin{longtable}{| p{.40\textwidth} | p{.50\textwidth} |}
\hline

\textbf{Nome do Caso de Uso} &
UC04 -- Manter Humor \\
\hline

\textbf{Ator Principal} &
Usuário \\
\hline

\textbf{Atores Secundários} &
--- \\
\hline

\textbf{Resumo} &
Permite registrar, editar e remover registros de humor. \\
\hline

\textbf{Pré-condições} &
Usuário autenticado. \\
\hline

\textbf{Pós-condições} &
O registro de humor é armazenado e passa a compor o histórico do usuário. \\
\hline

\multicolumn{2}{|c|}{\textbf{Cenário Principal}}\\
\hline

\textbf{Ações do Ator} &
\textbf{Resposta do Sistema} \\
\hline

1. Seleciona registrar humor. &
2. O sistema apresenta o formulário de registro. \\
\cline{1-2}

3. Informa o humor principal. &
\\
\cline{1-2}

4. Informa os níveis emocionais e demais informações. &
\\
\cline{1-2}

5. Adiciona observações, caso desejar. &
\\
\cline{1-2}

6. Confirma o registro. & \\
\cline{1-2}

&
7. O sistema valida os dados informados. \\
\cline{1-2}

&
8. O sistema registra as informações. \\
\cline{1-2}

&
9. Sistema agenda a atualização dos indicadores e da correlação
entre sono e energia. \\
\hline

\multirow{3}{.40\textwidth}{\textbf{Restrições e validações}}
&
O humor principal é obrigatório. \\
\cline{2-2}

&
Os níveis emocionais devem estar dentro do intervalo permitido. \\
\cline{2-2}

&
A data do registro é obrigatória. \\
\hline

\multicolumn{2}{|c|}{\textbf{Cenário Alternativo A1 -- Vincular gatilho}}\\
\hline
\textbf{Ações do Ator} &
\textbf{Resposta do Sistema} \\
\hline

1. Seleciona um ou mais gatilhos. & \\
\cline{1-2}

&
2. O sistema associa os gatilhos ao registro de humor. \\
\hline

\end{longtable}
